\documentclass[11pt]{article}

\usepackage[T1]{fontenc}
\usepackage{amsmath,amssymb,amsthm}
\usepackage[margin=1in]{geometry}
\usepackage[hidelinks]{hyperref}\hypersetup{pdftitle={Exact Counting Formulas for Two-Black-Square Crossword Rook Placements},pdfauthor={Xiaojian Zeng},pdfsubject={Combinatorics},pdfkeywords={crossword grids, rook placements, perfect matchings, permanents}}
\usepackage{microtype}

\newtheorem{theorem}{Theorem}
\newtheorem{lemma}[theorem]{Lemma}
\newcommand{\RP}{\operatorname{RP}}
\newcommand{\permanent}{\operatorname{per}}

\title{Exact Counting Formulas for Two-Black-Square\\
Crossword Rook Placements}
\author{Xiaojian Zeng\\
PBC School of Finance, Tsinghua University, 100190\\
\texttt{xiaojianzeng@foxmail.com}}
\date{}

\begin{document}
\maketitle

\begin{abstract}
For an $n\times n$ crossword grid with exactly two black squares, we give
coordinate-dependent factorial formulas for the number of complete
non-attacking rook placements in every relative-position class.  The formulas
distinguish the two order types of nonattacking interior squares and reduce the
fixed-two-black-square extremal problem of Lewis and Won to an explicit integer
optimization.  They also show that at $n=15$ a central edge-cross configuration
has $2432\cdot12!$ placements, exceeding the opposite-corner value
$2379\cdot12!$.  We do not determine the maximum for every $n$ or classify all
of its equality cases; the original extremal problem remains unresolved.
\end{abstract}

\section{Setting and contribution}

An \emph{across word} in a crossword grid is a maximal nonempty horizontal run
of white squares, and a \emph{down word} is defined vertically.  A complete
non-attacking rook placement meets every across and every down word exactly
once.  These conventions include words of length one and impose none of the
connectivity, symmetry, or minimum-word-length conditions of American
crosswords.

Lewis and Won introduced this model and proved that complete placements are
perfect matchings of the bipartite word-intersection graph
\cite[Section~2]{2609.03081}.  Thus, if $A_G$ is its $0$--$1$ biadjacency
matrix, then $|\RP(G)|=\permanent(A_G)$; see also
\cite[Chapter~7]{BrualdiRyser1991} for permanent background.  Lewis and Won
count the grids with zero or one black square and ask for the maximizing
arrangement at each fixed number of black squares
\cite[Question~2.2]{2609.03081}.  The interval structure of the present
$0$--$1$ matrices permits the exact matching decomposition below.

Fix $n\geq3$ and an unordered pair $B$ of black squares.  Write $G_B$ for the
resulting grid and
\[
   N(B)=|\RP(G_B)|.
\]
Coordinates are one-based.  A row or column coordinate is \emph{interior} if
it lies in $\{2,\ldots,n-1\}$.  We adopt the convention
\[
       m! =0 \qquad\text{when }m<0                                      \tag{1}
\]
only in displayed sums where such a term represents an impossible matching
case.

For integers $a,c\geq1$ and $b\geq0$, define
\begin{align*}
 D(a,b,c)&=ac,\\
 U(a,b,c)&=ac(a+b-1),\\
 V(a,b,c)&=ac(b+c-1),\\
 Q(a,b,c)&=ac\bigl(ab+ac-a+b^2+bc-3b-c+1\bigr).             \tag{2}
\end{align*}
If $x<y$ are interior coordinates, their \emph{profile} is
\[
       \pi_n(x,y)=(x-1,y-x-1,n-y).                           \tag{3}
\]
We append subscripts $r$ and $c$ to $D,U,V,Q$ when (2) is evaluated on
row and column profiles, respectively.

\section{The exhaustive formula}

\begin{theorem}[Two-black-square counting classification]\label{thm:main}
For every $n\geq3$ and every unordered two-square set $B$, the number $N(B)$
is given by exactly one of the following cases, together with their reflections
and transposes.  Every configuration not listed has $N(B)=0$.

\begin{enumerate}
\item If $B$ consists of opposite corners, then
\[
       N(B)=n!-2(n-1)!+(n-2)!.                               \tag{O}
\]

\item If $B=\{(1,1),(r,c)\}$ with $2\leq r,c\leq n-1$, then,
for $n\geq4$,
\[
 N(B)=(n-c)(n-r)(n-4)!
       \bigl((c-1)(r-1)(n-3)-(c-2)(r-2)\bigr).               \tag{C}
\]
For $n=3$, the unique such pair has $N(B)=1$.

\item If $B=\{(1,c),(r,1)\}$ with $2\leq r,c\leq n-1$, then
\[
 N(B)=(n-c)(n-r)
       \bigl((c-2)(r-2)(n-3)!+(n-2)!\bigr).                 \tag{E}
\]

\item Suppose both squares are interior and lie in distinct rows and
columns.  Let $r_-<r_+$ and $c_-<c_+$ be their row and column coordinates,
and evaluate (2) on $\pi_n(r_-,r_+)$ and $\pi_n(c_-,c_+)$.  Then
\[
 N(B)=Q_cQ_r(n-6)!+H(n-5)!+D_cD_r(n-4)!,                    \tag{I}
\]
where
\[
H=\begin{cases}
 U_cV_r+V_cU_r,&B=\{(r_-,c_-),(r_+,c_+)\}\quad
                  \text{(parallel order)},\\
 U_cU_r+V_cV_r,&B=\{(r_-,c_+),(r_+,c_-)\}\quad
                  \text{(crossed order)}.
\end{cases}                                                  \tag{4}
\]

\item If both squares lie in one boundary row and their column pair is one
of $\{1,2\}$, $\{1,n\}$, or $\{n-1,n\}$, then
\[
       N(B)=n!-2(n-1)!.                                      \tag{B}
\]
The transposed formula holds in a boundary column.

\item If $B=\{(r,c_1),(r,c_2)\}$, where
$2\leq r,c_1,c_2\leq n-1$ and $c_1+1<c_2$, then
\begin{align*}
N(B)={}&(c_1-1)(c_2-c_1-1)(n-c_2)\\
 &\mathrel{}\cdot(r-1)(r-2)(n-r)(n-r-1)(n-5)!.              \tag{S}
\end{align*}
The transposed formula holds for a common interior column.
\end{enumerate}
\end{theorem}

The theorem is a counting classification, not an extremal classification.
In particular, formula (I) leaves a nontrivial discrete optimization in the
row and column profiles.

\section{Proof}

We first show that the list of cases is exhaustive.

\begin{lemma}[Word-count reduction]\label{lem:words}
If the two black squares lie in distinct rows and columns, put
\[
 \alpha=\#\{\text{black squares in interior columns}\},\qquad
 \delta=\#\{\text{black squares in interior rows}\}.
\]
The grid has $n+\alpha$ across words and $n+\delta$ down words.  Consequently
$N(B)=0$ unless $\alpha=\delta$.

If the squares share a row and have columns $c_1<c_2$, put
\[
s={\bf1}_{c_1>1}+{\bf1}_{c_2>c_1+1}+{\bf1}_{c_2<n}.          \tag{5}
\]
There are $n-1+s$ across words.  There are $n$ down words when the common row
is a boundary row, and $n+2$ when it is interior.
\end{lemma}

\begin{proof}
In a row containing one black square, the white cells form two words exactly
when the black square is in an interior column; otherwise they form one word.
Distinct black rows therefore give $n+\alpha$ across words.  The column
argument is identical.  If the squares share a row, the other $n-1$ rows each
give one across word, while the affected row has the number of nonempty white
segments in (5).  Since the black squares occupy distinct columns, each of
those columns is split into two down words precisely when the common row is
interior.  This proves all assertions.
\end{proof}

A perfect matching requires equally many vertices in the two parts.  In the
distinct-row-and-column setting, $\alpha=\delta=0$ forces opposite corners.
When both equal one, either the same black square has its two coordinates
interior, giving a corner--interior pair, or the interior row and interior
column occur at different squares, giving an edge-cross pair.  When both equal
two, both squares are interior.  In a common boundary row, equality requires
$s=1$, whose solutions are precisely the three column pairs in case~(B).  In
a common interior row, equality requires $s=3$, which is exactly the hypothesis
of (S).  Transposition completes the exhaustion argument.

We now count each surviving class.  The recurring device is to order the split
interval words before the ordinary words in the word-biadjacency matrix.  Once
a matching among split words is fixed, every unmatched split across word must
be injected into an ordinary column, and every unmatched split down word into
an ordinary row.  The unused ordinary rows and columns form an all-one square
block, contributing a factorial.  Partitioning matchings by the chosen
split-to-split edges makes this decomposition disjoint and exhaustive.

For opposite corners, the word graph is $K_{n,n}$ with two vertex-disjoint
edges deleted.  Inclusion--exclusion gives (O).  In case (B), the two deleted
edges share a vertex, so no matching uses both; this gives (B).

For (C), first ignore the forbidden edge created by the corner square.  The
left and right across intervals at $(r,c)$ have respectively $c-1$ and $n-c$
ordinary-column choices, while the upper and lower down intervals have
$r-1$ and $n-r$ ordinary-row choices.  The remaining all-one core contributes
$(n-3)!$, giving
\[
 (c-1)(n-c)(r-1)(n-r)(n-3)!.
\]
The matchings counted here that use the forbidden corner edge are obtained by
first reserving that edge.  The four interval choices then contribute
$(c-2)(n-c)(r-2)(n-r)$ and the core contributes $(n-4)!$.
Subtracting and factoring proves (C).  Direct inspection gives the stated
$n=3$ value.

For (E), reflect so that the black squares are $(1,c)$ and $(r,1)$.  The sole
split-to-split edge is the white corner $(1,1)$.  If it is not used, the four
interval injections and the ordinary core contribute
\[
 (n-c)(n-r)(c-2)(r-2)(n-3)!.
\]
If it is used, the two remaining outer intervals have $(n-c)(n-r)$ choices
and the ordinary core contributes $(n-2)!$.  Their sum is (E).

For (S), the three horizontal intervals have disjoint sets of available
ordinary columns, of respective sizes $c_1-1$, $c_2-c_1-1$, and $n-c_2$.
The two upper vertical intervals require an ordered pair of distinct ordinary
rows above $r$, giving $(r-1)(r-2)$ choices.  The lower pair similarly gives
$(n-r)(n-r-1)$ choices.  The remaining core has size $n-5$, proving (S).

It remains to establish (I).  Consider two split coordinates $x<y$, and set
$(a,b,c)=\pi_n(x,y)$.  To inject all four interval words into ordinary
coordinates, the two outer intervals first force choices in disjoint regions
of sizes $a$ and $c$, giving $ac$ choices.  Conditional on this pair, the
other two choices contribute
\[
     b(a+b-2)+(c-1)(a+b-1).                                  \tag{6}
\]
Expanding $ac$ times (6) gives $Q(a,b,c)$.  If respectively one or the other
middle-facing interval word is removed, the same injection count is
$U(a,b,c)$ or $V(a,b,c)$; if both are removed it is $D(a,b,c)$.

For two black squares in distinct interior rows and columns, there are exactly
two white cross-intersections between split across and split down words.  They
are vertex-disjoint.  Selecting neither gives $Q_cQ_r(n-6)!$, and selecting
both gives $D_cD_r(n-4)!$.  Selecting exactly one removes one middle-facing
word on each side.  Tracking their orientations gives $U_cV_r+V_cU_r$ in
parallel order and $U_cU_r+V_cV_r$ in crossed order.  The remaining core has
size $n-5$.  This proves (I), and hence Theorem~\ref{thm:main}.

\section{A concrete extremal obstruction}

The opposite-corner formula agrees with the maximum in the small census
described below, but it is not maximal for all $n$.  At $n=15$, take
\[
       B=\{(1,7),(7,1)\}.
\]
Formula (E) gives
\[
 N(B)=8^2\bigl(5^2+13\bigr)12!=2432\cdot12!.
\]
Opposite corners give, by (O),
\[
 15!-2\cdot14!+13!=2379\cdot12!.
\]
Thus the edge-cross pair is larger by
$53\cdot12!=25{,}387{,}084{,}800$.  This refutes only the provisional
opposite-corner extrapolation; it neither supplies the global maximum nor
classifies equality.

\section{Exact regression evidence and limitations}

Two independent exact counters were run on every unordered pair for
$3\leq n\leq8$.  The first constructs the word graph and evaluates its
permanent by row-ordered subset dynamic programming.  The second uses direct
deletion recursion, branching on a minimum-degree across vertex.  A separate
symmetry routine forms all eight images generated by horizontal reflection,
vertical reflection, and transposition.  The two counters agree on all
$4{,}278$ configurations, represented by $620$ symmetry orbits, and the orbit
sizes recover the full labeled totals.  A separately coded evaluator of
(O)--(S) agrees with every orbit record.  All arithmetic is integral; there is
no floating-point or randomized step.  This architecture---symmetry-reduced
exhaustion, reconstruction of the full action, and checking by a distinct
exact implementation---also appears in finite combinatorial certification
outside rook theory \cite[Sections~3, 4, and 6]{202608.1658}.

This regression is evidence for the implementation, not a proof of the
all-$n$ theorem; the proof is the case decomposition above.  It also stops at
$n=8$, whereas the original extremal project requests a two-method census
through $n=12$.  Most importantly, the discrete maximization of (I) and the
comparison with the other classes remain undone.  Therefore this note does not
determine $M_2(n)$ for all $n$, and it does not resolve Question~2.2 even in its
full $k=2$ extremal form.

\section*{Computational and AI assistance disclosure}

OpenAI Codex was used to assist the organization of the case split, symbolic
formula development, construction of the exact Python checks, and preparation
of this manuscript.  The quantified mathematical argument is written out in
the text.  The companion Python programs and JSON census permit deterministic
recomputation of every finite check reported above.  Those computations do not
replace any all-$n$ proof step.

\bibliographystyle{plain}
\bibliography{references}

\end{document}
